\documentclass[11pt,a4paper]{article}

\usepackage[utf8]{inputenc}
\usepackage[T1]{fontenc}
\usepackage{amsmath,amsfonts,amssymb}
\usepackage{geometry}
\usepackage{enumitem}
\usepackage{titlesec}
\usepackage{tikz}

% Page Geometry
\geometry{top=1in, bottom=1in, left=1in, right=1in}

% Typography & Styling
\titleformat{\section}{\large\bfseries}{\thesection}{1em}{}[\titlerule]
\titleformat{\subsection}{\normalsize\bfseries}{\thesubsection}{1em}{}
\setlist[itemize]{leftmargin=*, noitemsep, topsep=2pt}

% Metadata
\title{\textbf{Psychological \& Intellectual Portrait:\\Maksym Sokhatskyi (Namdak Tonpa)}}
\author{Analytical Profile Note}
\date{\today}

\begin{document}

\maketitle

\begin{abstract}
This brief note outlines the cognitive, philosophical, and professional archetype of Ukrainian computer scientist and researcher Maksym Sokhatskyi. Synthesized from public contributions, research directions, and systemic architecture designs, he is characterized as \textbf{The Avant-Garde Rationalist}.
\end{abstract}

\section{Core Psychological Archetype}
Sokhatskyi represents a non-linear archetype that seamlessly blends the rigid, uncompromising logic of formal mathematics with the expansive, interconnected framework of Eastern philosophy. He stands at the intersection of ultra-formalism and philosophical rebellion.

\begin{figure}[h!]
\centering
\begin{tikzpicture}[node distance=2cm, auto, >=stealth, thick]
    % Styles
    \tikzstyle{block} = [rectangle, draw, text width=5cm, text centered, rounded corners, minimum height=1cm, fill=gray!10]
    \tikzstyle{core} = [rectangle, draw, text width=6cm, text centered, bold font, minimum height=1.2cm, fill=gray!20]
    
    % Nodes
    \node [core] (center) {THE AVANT-GARDE RATIONALIST};
    \node [block, below left=1cm and 0.5cm of center] (left) {\textbf{Ultra-Formalist Mind} \\ $\bullet$ Mathematical proofs \\ $\bullet$ Homotopy Type Theory \\ $\bullet$ Rejection of legacy clutter};
    \node [block, below right=1cm and 0.5cm of center] (right) {\textbf{Philosophical Rebel} \\ $\bullet$ Tibetan Buddhist worldview \\ $\bullet$ Anti-corruption tech drive \\ $\bullet$ Radical self-sovereignty};
    
    % Arrows
    \draw[->] (center) -- (left);
    \draw[->] (center) -- (right);
\end{tikzpicture}
\caption{Conceptual map of Sokhatskyi's intellectual psyche.}
\end{figure}

\section{Key Cognitive Pillars}

\subsection{The Ultra-Formalist Mind}
Sokhatskyi exhibits a cognitive profile that craves absolute certainty, structural elegance, and mathematical purity.
\begin{itemize}
    \item \textbf{Obsession with Formal Verification:} His dedication to Homotopy Type Theory (HoTT) and building automated provers (e.g., \textit{Anders}) shows a mind unsatisfied with ``good enough'' software. He views programming as a discipline of absolute logic where code must be mathematically proven correct.
    \item \textbf{Hyper-Reductionism:} His work with minimalist Erlang runtime structures indicates a cognitive preference for stripping away architectural bloat to find core, universal truths in systems design.
\end{itemize}

\subsection{The Non-Conformist Maverick}
While deeply analytical, he completely diverges from the traditional, institutional academic archetype.
\begin{itemize}
    \item \textbf{Anti-Corruption Ethos:} He actively positions himself against what he terms ``technological corruption'' in IT. This points to a highly principled, uncompromising moral framework. He is fiercely independent and vocal against legacy systems and bureaucratic engineering inefficiencies.
    \item \textbf{Radical Self-Sovereignty:} He heavily identifies with his digital and spiritual persona, \textit{Namdak Tonpa}, maintaining a distinct online identity and visual branding. He operates as an outsider-insider: capable of building massive national infrastructures while remaining culturally detached from mainstream corporate culture.
\end{itemize}

\subsection{The Union of Logic and Mysticism}
The most unique aspect of his psyche is the deliberate synthesis of two seemingly opposite worlds:
\begin{itemize}
    \item \textbf{Eastern Spiritual Framework:} As a practicing Buddhist studying Tibetan, his mental models are grounded in deep introspection, mindfulness, and the structural interdependence of reality.
    \item \textbf{Constructive Mathematics:} For Sokhatskyi, Buddhist philosophy and advanced type theory serve as extensions of the same pursuit: understanding the fundamental nature of reality, language, and thought through clean, unclouded observation.
\end{itemize}

\section{Summary Profile}
Sokhatskyi is an \textbf{Intellectual Purist}. He is driven by a profound disdain for systemic messiness—whether that messiness is a buggy piece of enterprise code, bureaucratic IT spending, or a cluttered mind. He finds peace and purpose in building frameworks that are unassailable, elegant, and perfectly structured from the ground up.

\end{document}

